% cuadro C12_2
\begin{table}[htbp]
\centering
\scriptsize
\caption{Composición del anexo de igualdad entre mujeres y hombres: los doce programas mayores}
\label{cua:C12_2}
\begin{tabular}{p{7.6cm}p{1.4cm}rr}
\toprule
\textbf{Programa} & \textbf{Clave} & \textbf{Millones} & \textbf{\% del anexo} \\
\midrule
Pensión para el Bienestar de las Personas Adultas Mayores & S176 & 261 776,1 & 43,7 \\
Beca Universal de Educación Básica Rita Cetina & S072 & 103 508,8 & 17,3 \\
La Escuela es Nuestra & S282 & 25 488,6 & 4,3 \\
Beca Universal de Educación Media Superior Benito Juárez & S311 & 21 279,5 & 3,6 \\
Programa de Vivienda Social & S177 & 19 968,0 & 3,3 \\
Servicios de guardería & E070 & 17 062,0 & 2,8 \\
Servicios de atención a la salud & E031 & 14 993,2 & 2,5 \\
Jóvenes Construyendo el Futuro & S280 & 14 975,4 & 2,5 \\
Servicios de Educación Superior y Posgrado & E023 & 14 940,0 & 2,5 \\
Pensión Mujeres Bienestar & S316 & 14 550,0 & 2,4 \\
Pensión para el Bienestar de las Personas con Discapacidad Permanente & S286 & 11 871,3 & 2,0 \\
Sembrando Vida & S287 & 7 820,0 & 1,3 \\
\bottomrule
\end{tabular}
\\[2pt]\begin{minipage}{0.92\textwidth}\scriptsize El 48.1 por ciento del anexo son programas de pensión. Fuente: base de anexos transversales del PPEF 2026. Elaborado por el ITED.\end{minipage}
\end{table}
